\everymath{\displaystyle}

%%%%%%%%%%%%%%%%%%%%%%%%%%%%%%%%%%%%%%%%
%           main body
%%%%%%%%%%%%%%%%%%%%%%%%%%%%%%%%%%%%%%%%

%-----------------------------------
%	TITLE PAGE
%-----------------------------------

\title[第九章\\
数项级数]{\texorpdfstring{\LARGE \bfseries 第九章\quad 数项级数}{第九章 数项级数}} % The short title appears at the bottom of every slide, the full title is only on the title page
\author[\smaller \S 9.4\\
任意项级数] % Your institution as it will appear on the bottom of every slide, maybe shorthand to save space
{\Large \bfseries \S 9.4 任意项级数}
% \author{Singular Value Decomposition} % Your name
% \date{\today} % Date, can be changed to a custom date
\date{}

%------------------------------------------------

\begin{document}

%------------------------------------------------

\begin{frame}

\titlepage % Print the title page as the first slide

\end{frame}

%-----------------------------------
%	PRESENTATION SLIDES
%-----------------------------------

\begin{frame}{从正项级数到任意项级数}

\begin{columns}
\begin{column}{0.55\textwidth}
一个级数, 如果只有有限个负项或有限个正项, 都可以用正项级数的~{\color{red}各种判别法}来判断它的敛散性.

\vspace{0.6em}

\only<1-6>{\phantom{如果一个级数即有无限个正项, 又有无限个负项, 那么正项级数的各种判别法不再适用!}}
\only<7->{如果一个级数即有无限个正项, 又有无限个负项, 那么正项级数的各种判别法不再适用!}
\end{column}

\begin{column}{0.42\textwidth}
\only<2->{
\begin{empheq}[box={\fancymath[colback=cyan!30, sharp corners]}]{align*}
& \only<2->{\text{\color{red!80}\ding{43} 比较判别法}} \\
& \only<3->{\text{\color{red!80}\ding{43} Cauchy 判别法}} \\
& \only<4->{\text{\color{red!80}\ding{43} d'Alembert 判别法}} \\
& \only<5->{\text{\color{red!80}\ding{43} Raabe 判别法}} \\
& \only<6->{\text{\color{red!80}\ding{43} 积分判别法}} \\
& \only<6->{\text{\color{red!80}\ding{43} $\cdots\cdots$}}
\end{empheq}
}
\end{column}
\end{columns}

\vspace{1em}

\only<8->{
\begin{center}
我们只能从最基本的~{\color{red}Cauchy 收敛原理}出发, 进行研究.
\end{center}
}

\end{frame}

%------------------------------------------------

\section[Cauchy\\收敛原理]{Cauchy 收敛原理}

%------------------------------------------------

\begin{frame}{级数的 Cauchy 收敛原理}

\begin{thm}[Cauchy 收敛原理]
级数 $\sum\limits_{n=1}^\infty x_n$ 收敛的充分必要条件是: 对任意给定的 $\varepsilon > 0,$ 存在正整数 $N,$ 使得
\[\left\lvert x_{n+1} + x_{n+2} + \cdots + x_m \right\rvert = \left\lvert \sum\limits_{k=n+1}^m x_k \right\rvert < \varepsilon\]
对一切 $m > n > N$ 成立.
\end{thm}

\pause
\vspace{0.6em}

\begin{attnblock}
\begin{itemize}
\item[\ding{43}] Cauchy 收敛原理对一般的度量空间都是成立的
\pause
\item[\ding{43}] 取 $m = n + 1,$ 我们即得到级数收敛的必要条件$\lim\limits_{n\to\infty} x_n = 0.$
\end{itemize}
\end{attnblock}

\end{frame}

%------------------------------------------------

\section[Leibniz\\级数]{Leibniz 级数}

%------------------------------------------------

\begin{frame}{Leibniz 级数}

\begin{Def}[交错级数]
如果级数 $\sum\limits_{n=1}^\infty x_n = \sum\limits_{n=1}^\infty (-1)^{n+1} u_n ~ (u_n > 0),$ 则称其为\textbf{交错级数}.
\end{Def}

\pause

\begin{Def}[Leibniz 级数]
如果级数 $\sum\limits_{n=1}^\infty x_n = \sum\limits_{n=1}^\infty (-1)^{n+1} u_n ~ (u_n > 0)$ 为交错级数, 并且数列 $\{u_n\}$ 单调减少且收敛于 $0,$ 则称此级数为\textbf{Leibniz 级数}.
\end{Def}

\pause

\begin{block}{Leibniz 级数的一些例子}
\vspace{-1.5em}
\[\sum\limits_{n=1}^\infty \dfrac{(-1)^{n+1}}{n^s} (s > 0), ~~~ \sum\limits_{n=2}^\infty (-1)^{n} \dfrac{\ln{n}}{n}, ~~~ \sum\limits_{n=1}^\infty (-1)^{n+1} \dfrac{n^2}{n^3 + 1}, ~~ \cdots\]
\end{block}

\end{frame}

%------------------------------------------------

\begin{frame}{Leibniz 级数}

有一些 Leibniz 级数, 则不那么明显能一眼看出:

\pause
\vspace{0.5em}

考虑级数 $\sum\limits_{n=1}^\infty \sin \left( \sqrt{n^2+1} \pi \right),$ 有
\begin{align*}
\sin \left( \sqrt{n^2+1} \pi \right) & = {\color{red} (-1)^n} \sin \left( \sqrt{n^2+1} \pi ~ {\color{red} - n\pi} \right) \\
& = (-1)^n \sin \left( \sqrt{n^2+1} - n \right) \pi \\
& = (-1)^n ~ \underbrace{\color{red} \sin \dfrac{\pi}{\sqrt{n^2+1} + n}}_{u_n}.
\end{align*}

\pause

容易看出 $\{u_n = \sin \dfrac{\pi}{\sqrt{n^2+1} + n}\}$ 是单调减小数列, 且
\[\lim\limits_{n\to\infty} \sin \dfrac{\pi}{\sqrt{n^2+1} + n} = 0,\]
可知 $\sum\limits_{n=1}^\infty \sin \left( \sqrt{n^2+1} \pi \right)$ 是 Leibniz 级数.

\end{frame}

%------------------------------------------------

\begin{frame}{Leibniz 判别法}

\begin{thm}[Leibniz 判别法]
\begin{center}
Leibniz 级数必定收敛.
\end{center}
\end{thm}

\pause
\vspace{0.5em}

\startproof[bg=yellow, fg=red]
首先, 记 $p = m - n,$ 有
\[\left\lvert x_{n+1} + x_{n+2} + \cdots + x_{n+p} \right\rvert = \left\lvert u_{n+1} - u_{n+2} + \cdots + (-1)^{p+1} u_{n+p} \right\rvert.\]
\pause
当 $p$ 是奇数时, 有
\begin{multline*}
u_{n+1} - u_{n+2} + \cdots + (-1)^{p+1} u_{n+p} \\
= \begin{cases}
(u_{n+1} - u_{n+2}) + \cdots + (u_{n+p-2} - u_{n+p-1}) + u_{n+p} ~ {\color{red} \geqslant u_{n+p}}, \\
u_{n+1} - (u_{n+2} - u_{n+3}) - \cdots - (u_{n+p-1} - u_{n+p}) ~ {\color{red} \leqslant u_{n+1}}
\end{cases}
\end{multline*}

\end{frame}

%------------------------------------------------

\begin{frame}{Leibniz 判别法的证明}

当 $p$ 是偶数时, 有
\[
u_{n+1} - u_{n+2} + \cdots + (-1)^{p+1} u_{n+p}
= \begin{cases}
(u_{n+1} - u_{n+2}) + \cdots + (u_{n+p-2} - u_{n+p-1}) ~ {\color{red} \geqslant 0}, \\
u_{n+1} - (u_{n+2} - u_{n+3}) - \cdots - u_{n+p} ~ {\color{red} \leqslant u_{n+1}}.
\end{cases}
\]

\pause

总之, 成立
\begin{align*}
\left\lvert x_{n+1} + x_{n+2} + \cdots + x_{n+p} \right\rvert & = \left\lvert u_{n+1} - u_{n+2} + \cdots + (-1)^{p+1} u_{n+p} \right\rvert \\
& \leqslant u_{n+1}.
\end{align*}

由于 $\lim\limits_{n\to\infty} u_n = 0,$ 所以对于任意给定的 $\varepsilon > 0,$ 存在正整数 $N,$ 使得对一切 $n > N,$ 成立
\only<2>{
\[\varepsilon > u_{n+1},\]
}
\only<3->{
\[\varepsilon > u_{n+1} ~ {\color{red} \geqslant \left\lvert x_{n+1} + x_{n+2} + \cdots + x_{n+p} \right\rvert},\]
}
\only<4->{
由 Cauchy 收敛原理知 Leibniz 级数必收敛.
\qed
}

\end{frame}

%------------------------------------------------

\section[A-D 判别法]{Abel 判别法与 Dirichlet 判别法}

%------------------------------------------------

\begin{frame}{回忆 | {\smaller[2] 反常积分的A-D判别法}}

我们先回忆一下一般函数的反常积分敛散性的 Abel 判别法与 Dirichlet 判别法 (A-D 判别法):

\pause

考虑 $\int_a^{+\infty} f(x)g(x) \mathrm{d} x,$ 有
\begin{itemize}
\item[(1)] Abel 判别法:
\[\int_a^{+\infty} f(x) \mathrm{d} x ~ \text{收敛}, ~ g(x) ~ \text{单调有界} \implies \int_a^{+\infty} f(x)g(x) \mathrm{d} x ~ \text{收敛};\]
\item[(2)] Dirichlet 判别法:
\[\int_a^{A} f(x) \mathrm{d} x ~ \text{关于$A$有界}, ~ g(x) ~ \text{单调}, ~ \lim\limits_{x\to\infty} g(x) = 0 \implies \int_a^{+\infty} f(x)g(x) \mathrm{d} x ~ \text{收敛}\]
\end{itemize}

\pause

接下来要考虑的一般级数, 也假设它具有 $\sum\limits_{n=1}^{\infty} a_n b_n$ 的形式, 我们有类似的结论, 证明方法也是类似的.

\end{frame}

%------------------------------------------------

\begin{frame}{级数的 A-D 判别法}

\begin{thm}[级数的 A-D 判别法]
若下列条件之一满足, 则级数$\sum\limits_{n=1}^{\infty} a_n b_n$收敛:
\begin{itemize}
\item[(1)] Abel 判别法:
\[\sum\limits_{n=1}^{\infty} b_n ~ \text{收敛}, ~~~ a_n ~ \text{单调有界}\]
\item[(2)] Dirichlet 判别法:
\[\sum\limits_{n=1}^{N} b_n ~ \text{关于$N$有界}, ~~~ a_n ~ \text{单调趋于$0$}\]
\end{itemize}
\end{thm}

\end{frame}

%------------------------------------------------

\begin{frame}{A-D 判别法的分析与证明}

对于一般的情况, 目前能用的工具只有 Cauchy 收敛原理, 所以要考察
\[\left\lvert \sum\limits_{k=n+1}^{m} a_k b_k \right\rvert.\]

\pause

(1). Abel 判别法: 条件为 $\sum\limits_{n=1}^{\infty} b_n$ 收敛, $a_n$单调有界.
\vspace{-1em}
\begin{align*}
\sum\limits_{n=1}^{\infty} b_n ~ \text{收敛} ~~ & ~ {\color{red}\implies} ~~
\only<2>{\phantom{\forall~\varepsilon > 0, \exists N, \text{ 使得 } \forall~m > n > N, \left\lvert \sum\limits_{k=n+1}^{m} b_k \right\rvert < \varepsilon.}}
\only<3->{\forall~\varepsilon > 0, \exists N, \text{ 使得 } \forall~m > n > N, \left\lvert \sum\limits_{k=n+1}^{m} b_k \right\rvert < \varepsilon.} \\
\only<4->{a_n ~ \text{单调有界} ~~ & ~ {\color{red}\implies}} ~~ \only<5->{\exists M > 0, \text{ 使得 } \lvert a_n \rvert < M \text{ 对所有$n$成立}}
\end{align*}

\vspace{-1em}

\only<6->{
\begin{empheq}[box={\fancymath[colback=cyan!30, drop lifted shadow, sharp corners]}]{equation*}
\text{\bfseries\color{red!80}希望: $\sum\limits_{k=n+1}^{m} a_k b_k \rightsquigarrow \left( \sum\limits_{k=n+1}^{m} b_k \right) \cdot \big( a_{n+1}, \cdots, a_m ~ \text{相关的项} \big)$}
\end{empheq}
}

\end{frame}

%------------------------------------------------

\begin{frame}{A-D 判别法的分析与证明}

事实上, 如果记 $B_t = \sum\limits_{k=n+1}^{t} b_k, t > n,$ 并约定~${\color{red}B_n = 0},$ 那么所有的 $B_t$ 都满足 $\lvert B_t \rvert < \varepsilon,$ 并且有
{\smaller[1]
\begin{align*}
\left\lvert \sum\limits_{k=n+1}^{m} a_k b_k \right\rvert & = \left\lvert \sum\limits_{k=n+1}^{m} a_k \left( B_{k} - B_{k-1} \right) \right\rvert = \left\lvert \sum\limits_{k=n+1}^{m} a_k B_{k} - \sum\limits_{k = ~ {\color{red}n+1}}^{{\color{red}m}} a_{{\color{red}k}} B_{{\color{red}k-1}} \right\rvert \\
& = \left\lvert \sum\limits_{k=n+1}^{m} a_k B_{k} - \sum\limits_{k = ~ {\color{red}n}}^{{\color{red}m-1}} a_{{\color{red}k+1}} B_{{\color{red}k}} \right\rvert \\
& = \left\lvert a_mB_m - \sum\limits_{k=n+1}^{m-1} (a_{k+1} - a_k) B_{k} + a_{n+1}{\color{red}B_n} \right\rvert \\
& \leqslant \left\lvert a_m \right\rvert \varepsilon + \sum\limits_{k=n+1}^{m-1} \left\lvert a_{k+1} - a_k \right\rvert \varepsilon
\end{align*}
}

\end{frame}

%------------------------------------------------

\begin{frame}{A-D 判别法的分析与证明}

由于 $a_n$ 单调, 所以 $\sum\limits_{k=n+1}^{m-1} \left\lvert a_{k+1} - a_k \right\rvert = \left\lvert a_{n+1} - a_{m} \right\rvert,$ 于是
\begin{align*}
\left\lvert \sum\limits_{k=n+1}^{m} a_k b_k \right\rvert & \leqslant \left\lvert a_m \right\rvert \varepsilon + {\color{red} \sum\limits_{k=n+1}^{m-1} \left\lvert a_{k+1} - a_k \right\rvert} \varepsilon \\
& = \left( \left\lvert a_m \right\rvert + {\color{red} \left\lvert a_{n+1} - a_{m} \right\rvert} \right) \varepsilon \\
& < 3M\varepsilon
\end{align*}
\pause
有 Cauchy 收敛原理, 阿贝判别法得证.

\end{frame}

%------------------------------------------------

\begin{frame}{A-D 判别法的分析与证明}

(2). Dirichlet 判别法: 条件为$\sum\limits_{n=1}^{N} b_n$ 关于 $N$ 有界, $a_n$ 单调趋于 $0.$
\begin{align*}
a_n ~ \text{单调趋于 $0$} ~~ & {\color{red}\implies} ~~
\only<1>{\phantom{\forall~\varepsilon > 0, \exists N, ~ \text{使得} ~ \forall~n > N, \left\lvert a_n \right\rvert < \varepsilon.}}
\only<2->{\forall~\varepsilon > 0, \exists N, ~ \text{使得} ~ \forall~n > N, \left\lvert a_n \right\rvert < \varepsilon.} \\
\only<3->{\sum\limits_{n=1}^{N} b_n ~ \text{关于 $N$ 有界} ~~ & {\color{red}\implies}} ~~ \only<4->{\exists M > 0, ~ \text{使得} ~ \forall~N ~ \text{有} ~ \left\lvert \sum\limits_{n=1}^{N} b_n \right\rvert < ~ {\color{red}M}}
\end{align*}

\only<5->{
利用证明阿贝判别法类似的方法, 有
\begin{align*}
\left\lvert \sum\limits_{k=n+1}^{m} a_k b_k \right\rvert & = \left\lvert a_m ~ {\color{red}B_m} - \sum\limits_{k=n+1}^{m-1} (a_{k+1} - a_k) ~ {\color{red}B_{k}} \right\rvert \\
& \leqslant \left\lvert a_m \right\rvert \cdot ~ {\color{red}2M} + \sum\limits_{k=n+1}^{m-1} \left\lvert a_{k+1} - a_k \right\rvert \cdot ~ {\color{red}2M} \\
& = \left( \left\lvert a_m \right\rvert + \left\lvert a_{n+1} - a_{m} \right\rvert \right) \cdot 2M \\
& < 6M\varepsilon
\end{align*}
}

\end{frame}

%------------------------------------------------

\begin{frame}{A-D 判别法的分析与证明}

反常积分敛散性的 A-D 判别法 (实际是证明过程中用到的积分第二中值定理) 与级数敛散性的 A-D 判别法的证明实际上都用到了如下的和式转换方法

\begin{lem}[Abel 变换 (分部求和公式)]
设 $\{a_n\}, \{b_n\}$ 是数列, 记 $B_n = \sum\limits_{k=1}^n b_k,$ $n = 1, 2, \dots,$ 那么
\[\sum\limits_{k=1}^n a_kb_k = a_n B_n - \sum\limits_{k=1}^{n-1} (a_{k+1} - a_k) B_k.\]
\end{lem}

\end{frame}

%------------------------------------------------

\begin{frame}{A-D 判别法的分析与证明}

相应地, 我们对 $\sum\limits_{k=1}^n a_kb_k$ 有如下的估计

\begin{lem}[Abel 引理]
设 $\{a_n\}$ 为单调数列, $\{B_n = \sum\limits_{k=1}^n b_k\}$ 为有界数列, $M_0$ 为它的一个界, 那么
\[\left\lvert \sum\limits_{k=1}^n a_kb_k \right\rvert \leqslant M_0 \cdot \left( \left\lvert a_n \right\rvert + \left\lvert a_n - a_1 \right\rvert \right).\]
\pause
特别地, 如果进一步有 $m \leqslant B_n \leqslant M,$ 并且 $\{a_n\}$ 非负不增, 那么
\[m a_1 \leqslant \sum\limits_{k=1}^n a_kb_k \leqslant M a_1.\]
\end{lem}

\end{frame}

%------------------------------------------------

\begin{frame}{A-D 判别法的几个注记}

1. 回过头来看 Leibniz 判别法:
\begin{center}
Leibniz 级数都收敛
\end{center}
\pause
\begin{empheq}[box={\fancymath[colback=cyan!30, drop lifted shadow, sharp corners]}]{equation*}
\text{\bfseries\color{red!80}Leibniz 级数: $\sum\limits_{n=1}^{\infty} (-1)^{n+1} u_n, ~~ 0 < u_n ~ \text{单调减小, 且} ~ u_n \to 0.$}
\end{empheq}

\pause
\vspace{0.6em}

若令 $b_n = (-1)^{n+1}, a_n = u_n,$ 那么它们满足 Dirichlet 判别法的条件
\begin{empheq}[box={\fancymath[colback=cyan!30, drop lifted shadow, sharp corners]}]{equation*}
\color{red!80} \sum\limits_{n=1}^{N} b_n ~ \text{关于 $N$ 有界}, ~~~ a_n ~ \text{单调趋于 $0$.}
\end{empheq}

\pause

也就是说, Leibniz 判别法可以由 Dirichlet 判别法导出, 可视为是 Dirichlet 判别法的一个特例.

\end{frame}

%------------------------------------------------

\begin{frame}{A-D 判别法的几个注记}

2. Abel 判别法可以由 Dirichlet 判别法导出:

\pause
\vspace{0.8em}

事实上, 若 Abel 判别法的条件成立: $\sum\limits_{n=1}^{\infty} b_n$ 收敛, $a_n$ 单调有界. \par
由数列的单调有界定理, $a_n$ 收敛于某个实数 $a = \lim\limits_{n\to\infty} a_n.$ 那么 $c_n : = a_n - a$ 就构成了一个单调趋于 $0$ 的数列.

\pause
\vspace{0.8em}

另一方面, 由于 $\sum\limits_{n=1}^{\infty} b_n$ 收敛, 那么 $\sum\limits_{n=1}^{N} b_n$ 显然关于 $N$ 有界. 于是, 根据 Dirichlet 判别法, $\sum\limits_{n=1}^{\infty} c_nb_n$ 收敛, 从而知
\[\sum\limits_{n=1}^{\infty} a_nb_n = \sum\limits_{n=1}^{\infty} c_nb_n + a \sum\limits_{n=1}^{\infty} b_n\]
也是收敛的.

\end{frame}

%------------------------------------------------

\begin{frame}{A-D判别法的应用}

\begin{eg}
设数列 $\{a_n\}$ 单调趋于 $0,$ 则 $\forall~x \in \mathbb{R},$ 级数 $\sum\limits_{n=1}^\infty a_n \sin nx$ 收敛.
\end{eg}

\pause
\vspace{0.4em}

\startproof[bg=yellow, fg=red]
$\{a_n\}$ 单调趋于 $0,$ 这一条件是 Dirichlet 判别法的一半; 自然希望证明另一半条件, 即 $\sum\limits_{n=1}^N \sin nx$ 关于 $N$ 有界.

\pause

\textbf{方法一:} 利用 $\sin nx = \dfrac{e^{inx} - e^{-inx}}{2i},$ 可以将 $\sum\limits_{n=1}^N \sin nx$ 表示为两个等比数列求和 ($x \neq 2k\pi$):
\begin{align*}
\sum\limits_{n=1}^N \sin nx & = \dfrac{1}{2i} \left( \sum\limits_{n=1}^N \left(e^{ix}\right)^n - \sum\limits_{n=1}^N \left(e^{-ix}\right)^n \right) \\
& = \dfrac{1}{2\sin\frac{x}{2}} \left( \cos\dfrac{x}{2} - \cos \left(N + \dfrac{1}{2}\right) x \right)
\end{align*}

\end{frame}

%------------------------------------------------

\begin{frame}{A-D 判别法的应用}

\textbf{方法二:} 利用三角函数的和差化积、积化和差公式:
{
% \smaller[1]
\begin{align*}
2\sum\limits_{n=1}^N \sin nx & = \sin x + \left( \sin x + \sin 2x \right) + \cdots + \sin Nx \\
& = \sin x + ~ {\color{red} \sin \dfrac{3}{2}x} \cos \dfrac{1}{2}x + \cdots + ~ {\color{red} \sin \dfrac{2N-1}{2}x} \cos \dfrac{1}{2}x + \sin Nx
\end{align*}
}
\pause
两边同时乘以 ${\color{red} \sin \dfrac{1}{2}x},$ 再利用
\[\sin \dfrac{1}{2}x \sin \sin \dfrac{2k-1}{2}x = \cos (k-1)x - \cos kx,\]
相邻项交错相消, 有
{
% \smaller[1]
\begin{align*}
2\sin \dfrac{1}{2}x\sum\limits_{n=1}^N \sin nx
= & \sin \dfrac{1}{2}x \sin x + \cos \dfrac{1}{2}x \cos x - \cos \dfrac{1}{2}x \cos Nx \\
& + \sin \dfrac{1}{2}x \sin Nx
= \cos\dfrac{x}{2} - \cos \left(N + \dfrac{1}{2}\right) x
\end{align*}
}

\end{frame}

%------------------------------------------------

\begin{frame}{A-D 判别法的应用}

总之, 当 $x \neq 2k\pi$ 时, 有
\[\sum\limits_{n=1}^N \sin nx = \dfrac{1}{2\sin\frac{x}{2}} \left( \cos\dfrac{x}{2} - \cos \left(N + \dfrac{1}{2}\right) x \right),\]
对于固定的 $x,$ 上式关于 $N$ 显然是有界的, $\dfrac{1}{\left\lvert \sin\frac{x}{2} \right\rvert}$ 即为它的一个界. 于是, 根据 Dirichlet 判别法, $\sum\limits_{n=1}^\infty a_n \sin nx$ 收敛.

\pause
\vspace{0.6em}

当 $x = 2k\pi$ 时, $\sum\limits_{n=1}^\infty a_n \sin nx$ 的每一项都是 $0,$ 显然也是收敛的.
\qed

\end{frame}

%------------------------------------------------

\section{级数的乘法}  % 未完成

%------------------------------------------------

\begin{frame}{级数的乘法}

\begin{question}
对两个 (收敛的) 级数 $\sum\limits_{n=1}^\infty a_n$ 和 $\sum\limits_{n=1}^\infty b_n,$ 如何定义 (计算) 它们的乘积
\[
\left( \sum\limits_{n=1}^\infty a_n \right) \cdot \left( \sum\limits_{n=1}^\infty b_n \right),
\]
是否还满足 \textbf{分配律}? 或者说, 分配律如何定义?
\[
\tikz[baseline=(m.north west)]{
\node (m) {
$\begin{aligned}
& a_1 b_1 + a_1 b_2 + a_1 b_3 + \cdots \\
+ & a_2 b_1 + a_2 b_2 + a_2 b_3 + \cdots \\
+ & a_3 b_1 + a_3 b_2 + a_3 b_3 + \cdots \\
& \vdots
\end{aligned}$
};

\draw[->, color=red] (m.north west) -- ++(5,0);
\draw[->, color=red] (m.north west) -- ++(0,-2.5);
}
\]

\pause

归根结底的问题是, 如何定义\textbf{部分和}?
\end{question}

\end{frame}

%------------------------------------------------

\begin{frame}{级数的乘法 -- 思路 1}

\begin{block}{思路 1}
分别取 $a_k$ 的前 $n_1$ 项以及 $b_k$ 的前 $n_2$ 项, 对应于上面那个图中左上角 $n_1 \times n_2$ 的矩形区域. 特别地, 将其中所有正方形 (即 $n_1 = n_2 = n$) 拿出来, 记
\[
d_n = \sum\limits_{1 \leqslant k, \ell \leqslant n} a_k b_\ell,
\]
并令 (形式地, 先不考虑 $d_n$ 是否收敛)
\[
\left( \sum\limits_{n=1}^\infty a_n \right) \cdot \left( \sum\limits_{n=1}^\infty b_n \right) = \lim\limits_{n\to\infty} d_n.
\]
\end{block}

\pause

下面, 我们来考虑这种思路的适用范围

\begin{prop}
如果 $\sum\limits_{n=1}^\infty a_n = A$ 和 $\sum\limits_{n=1}^\infty b_n = B$ 都是收敛的, 那么 $\lim\limits_{n\to\infty} d_n = AB.$
\end{prop}

\end{frame}

%------------------------------------------------

\begin{frame}{级数的乘法 -- 思路 1 的适用范围}

\startproof[bg=yellow, fg=red]
$\forall \varepsilon > 0,$ 存在正整数 $N,$ 使得对一切 $n > N,$ 成立
\[
\left\lvert \sum\limits_{k=1}^n a_k - A \right\rvert < \varepsilon / (|A| + |B| + 1), ~~ \text{且} ~~ \left\lvert \sum\limits_{k=1}^n b_k - B \right\rvert < \varepsilon / (|A| + |B| + 1).
\]
记 $\tau_1 = \sum\limits_{k=1}^n a_k -A , \tau_2 = \sum\limits_{k=1}^n b_k - B,$ 那么 $|\tau_1|, |\tau_2| < \varepsilon / (|A| + |B| + 1).$ 于是, 对一切 $n > N,$ 成立
\[
d_n = \left(  \sum\limits_{k=1}^n a_k \right) \left( \sum\limits_{k=1}^n b_k \right) = (A + \tau_1)(B + \tau_2) = AB + A\tau_2 + B\tau_1 + \tau_1\tau_2,
\]
那么有
\[
\left\lvert d_n - AB \right\rvert = \left\lvert A\tau_2 + B\tau_1 + \tau_1\tau_2 \right\rvert \leqslant (|A| + |B| + 1) \cdot \varepsilon / (|A| + |B| + 1) = \varepsilon.
\]

\end{frame}

%------------------------------------------------

\begin{frame}[fragile]{级数的乘法 -- 思路 2 (Cauchy 乘积)}

\begin{block}{思路 2 (Cauchy 乘积)}
按照反对角线方向, 将下标和为定值 $n+1$ 的项 $a_k b_\ell$ ($k + \ell = n + 1$) 归并在一起, 记
\[
c_n = \sum_{k=1}^{n} a_k b_{n-k+1} = a_1 b_n + a_2 b_{n-1} + \cdots + a_n b_1,
\]
并定义
\[
\left( \sum_{n=1}^\infty a_n \right) \cdot \left( \sum_{n=1}^\infty b_n \right) = \sum_{n=1}^\infty c_n.
\]
\end{block}

\[
\tikz[baseline=(m.north west)]{
\matrix (m) [matrix of math nodes,
            column sep={0.2cm, 0.1cm, 0.1cm, 0.1cm},
            row sep=0.1cm] {
    c_1 & a_1 b_1 & a_1 b_2 & a_1 b_3 & \cdots \\
    c_2 & a_2 b_1 & a_2 b_2 & a_2 b_3 & \cdots \\
    c_3 & a_3 b_1 & a_3 b_2 & a_3 b_3 & \cdots \\
    \vdots & \vdots & \vdots & \vdots & \ddots \\
};

% 第1条对角线: a_1b_1 (第1行第2列)
\draw[dashed, blue, thick] (m-1-2.north east) -- (m-1-2.south west);

% 第2条对角线: a_1b_2 (第1行第3列) 到 a_2b_1 (第2行第2列)
\draw[dashed, blue, thick] (m-1-3.north east) -- (m-2-2.south west);

% 第3条对角线: a_1b_3 (第1行第4列) 到 a_3b_1 (第3行第2列)
\draw[dashed, blue, thick] (m-1-4.north east) -- (m-3-2.south west);
}
\]

\end{frame}

%------------------------------------------------

\begin{frame}{级数的乘法 -- 思路 2 (Cauchy 乘积)}

Cauchy 乘积实际上是一种 \alert{卷积 (Convolution)} 的操作:
\begin{align*}
\text{连续版本:} & \quad f * g (x) = \int_{-\infty}^{+\infty} f(t) g(x - t) ~ \mathrm{d} t; \\
\text{离散有限版本:} & \quad A * B [k] = \sum_{i+j=k} A[i] \cdot B[j], ~~ \text{$A, B$ 的长度分别为 $n$ 和 $m$}; \\
\text{离散无限版本:} & \quad A * B [k] = \sum_{i+j=k} A[i] \cdot B[j].
\end{align*}

下面, 我们来研究这种思路下, 级数乘积的敛散性

\pause

\begin{thm}
设级数 $\sum\limits_{n=1}^\infty a_n = A$ 和 $\sum\limits_{n=1}^\infty b_n = B$ 都是收敛的, 且至少其中一个是绝对收敛的, 那么 Cauchy 乘积 $\sum\limits_{n=1}^\infty c_n$ 也是收敛的, 并且 $\sum\limits_{n=1}^\infty c_n = AB.$

\pause

若 $\sum\limits_{n=1}^\infty a_n, \sum\limits_{n=1}^\infty b_n$ 都只是条件收敛的, 那么 $\sum\limits_{n=1}^\infty c_n$ 可能发散.
\end{thm}

\end{frame}

%------------------------------------------------

\begin{frame}{Cauchy 乘积收敛性的分析与证明}

\startproof[bg=yellow, fg=red](分析)
我们设 $\sum\limits_{n=1}^\infty a_n = A$ 绝对收敛. 要证明 $\sum\limits_{n=1}^\infty c_n$ 收敛 (暂时不证明它的和是 $AB$), 则由 Cauchy 收敛准则, 需要考察
\[
\left| \sum_{n = n_1+1}^{n_2+1} c_n \right| = \left| \sum_{n = n_1+1}^{n_2+1} \sum_{i+j=n} a_i b_j \right|.
\]
\pause
上述和实际上对应无限矩阵中左上角的一块``梯形'':
\begin{columns}
\begin{column}{0.5\textwidth}
\begin{align*}
& \hspace{5em} a_1 {\color{red} b_{n_1}} + a_1 {\color{red} b_{n_1+1}} + \cdots + a_1 {\color{red} b_{n_2}} \\
& \hspace{3em} + a_2 {\color{red} b_{n_1-1}} + a_2 {\color{red} b_{n_1}} + \cdots + a_2 {\color{red} b_{n_2-1}} \\
& \hspace{1.5em} \reflectbox{$\ddots$} \hspace{11.5em} \reflectbox{$\ddots$} \\
& + {\color{red} a_{n_1}} b_1 + {\color{red} a_{n_1}} b_2 + \cdots + {\color{red} a_{n_1}} {\color{red} b_{n_2 - n_1 + 1}} \\
& + {\color{red} a_{n_1+1}} {\color{red} b_1} + \cdots + {\color{red} a_{n_1+1}} {\color{red} b_{n_2 - n_1}} \\
& \hspace{1.5em} \vdots \hspace{6.5em} \reflectbox{$\ddots$} \\
& + {\color{red} a_{n_2-1}} b_1 + {\color{red} a_{n_2-1}} b_2 \\
& + {\color{red} a_{n_2}} b_1
\end{align*}
\end{column}
\pause
\begin{column}{0.5\textwidth}
从左式可以看到, 对于 $\color{red} b_{n_1}, a_{n_2}, \dots$ 这些项, 当 $n_1, n_2$ 足够大时, 它们的绝对值可以被控制; 但是对于 $b_1, a_1, \dots$ 这些项, 它们的绝对值都无法被控制. 这就启发了我们, 需要将上式截断, 分而治之, 一部分用绝对值较小的 $a_n$ 来控制, 另一部分用绝对值较小的 $b_n$ 来控制.
\end{column}
\end{columns}

\end{frame}

%------------------------------------------------

\begin{frame}{Cauchy 乘积收敛性的分析与证明}

\begin{columns}
\begin{column}{0.5\textwidth}
\begin{align*}
\tikzcoord{L} & \hspace{5em} a_1 {\color{red} b_{n_1}} + a_1 {\color{red} b_{n_1+1}} + \cdots + a_1 {\color{red} b_{n_2}} \tikzcoord{R} \\
& \hspace{3em} + a_2 {\color{red} b_{n_1-1}} + a_2 {\color{red} b_{n_1}} + \cdots + a_2 {\color{red} b_{n_2-1}} \\
& \hspace{1.5em} \reflectbox{$\ddots$} \hspace{11.5em} \reflectbox{$\ddots$} \\
% {\only<2->{& \tikz[remember picture, overlay] \draw[dashed, blue] ([yshift=-3.2em]L) -- ([yshift=-3.2em]R); \\[-\baselineskip]}}
& + {\color{red} a_{n_1}} b_1 + {\color{red} a_{n_1}} b_2 + \cdots + {\color{red} a_{n_1}} {\color{red} b_{n_2 - n_1 + 1}} \\
& + {\color{red} a_{n_1+1}} {\color{red} b_1} + \cdots + {\color{red} a_{n_1+1}} {\color{red} b_{n_2 - n_1}} \\
& \hspace{1.5em} \vdots \hspace{6.5em} \reflectbox{$\ddots$} \\
& + {\color{red} a_{n_2-1}} b_1 + {\color{red} a_{n_2-1}} b_2 \\
& + {\color{red} a_{n_2}} b_1
\end{align*}
\end{column}
\pause
\begin{column}{0.5\textwidth}
记 $A_{i}^{j} = \sum\limits_{n=i}^{j} a_n, B_{i}^{j} = \sum\limits_{n=i}^{j} b_n,$ 那么左式以橫的视角来看, 可以表示为
\[
\sum a_* B_*^*,
\]
以竖的视角来看, 可以表示为
\[
\sum A_*^* b_*.
\]
\pause
由于 $\sum\limits_{n=1}^\infty a_n$ 绝对收敛, 我们使用第一种视角来控制上式的绝对值.
{\tikz[remember picture, overlay] \draw[dashed, red, thick] ([yshift=-3.2em]L) -- ([yshift=-3.2em]R);}
\end{column}
\end{columns}

\pause

记
\[
S := \sum\limits_{n=n_1+1}^{n_2+1} c_n = a_1 B_{n_1}^{n_2} + a_2 B_{n_1-1}^{n_2-1} + \cdots + a_{n+1} B_{n_1-n}^{n_2-n} + \cdots + a_{n_1} B_1^{n_2-n_1+1} + a_{n_1+1} B_1^{n_2-n_1} + \cdots + a_{n_2} B_1^1.
\]

\end{frame}

%------------------------------------------------

\begin{frame}{Cauchy 乘积收敛性的分析与证明}

由已知条件 $\sum\limits_{n=1}^\infty b_n = B$ 收敛, 那么
\[
\begin{cases}
\exists M > 0, ~ \text{使得} ~ \forall~N, ~ \left\lvert \sum\limits_{n=1}^N b_n \right\rvert < M ~ \implies \left\lvert B_i^j \right\rvert < 2M; \\
\forall~\varepsilon > 0, ~ \exists N_0, ~ \text{使得} ~ \forall~j > i > N_0, ~ \left\lvert B_i^j \right\rvert < \varepsilon.
\end{cases}
\]
\pause
由已知条件 $\sum\limits_{n=1}^\infty a_n = A$ 绝对收敛, 那么
\[
\begin{cases}
\exists M' > 0, ~ \text{使得} ~ \forall~N, ~ \sum\limits_{n=1}^N |a_n| < M'; \\
\forall~\varepsilon > 0, ~ \exists N_1, ~ \text{使得} ~ \forall~i > N_1, ~ |a_i| < \varepsilon.
\end{cases}
\]
\pause
取 $N = \max\{N_0, N_1\},$ $n_2, n_1 > 2N + 4,$ 并令 $n_0 = \lfloor n_1 / 2 \rfloor - 1 > N,$ 那么
\begin{align*}
|S| & \leqslant \left( |a_1| \cdot |B_{n_1}^{n_2}| + \cdots + |a_{n_0}| \cdot |B_{n_1-n_0+1}^{n_2-n_0+1}| \right) + \left( |a_{n_0+1}| \cdot |B_{n_1-n_0}^{n_2-n_0}| + \cdots + |a_{n_2}| \cdot |B_1^1| \right) \\
& < (|a_1| + \cdots + |a_{n_0}|) \cdot \varepsilon + \varepsilon \cdot 2M < (M' + 2M) \cdot \varepsilon.
\end{align*}

\end{frame}

%------------------------------------------------

\begin{frame}{Cauchy 乘积收敛性的分析与证明}

\begin{remark}
至于 $\sum\limits_{n=1}^\infty c_n = AB$ 的证明, 可以考察 $\left\lvert \sum\limits_{n=1}^N c_n - AB \right\rvert,$ 用类似的方法分析误差, 这里不再赘述.
\end{remark}

\pause

\begin{block}{$\sum\limits_{n=1}^\infty a_n,$ $\sum\limits_{n=1}^\infty b_n$ 都只是条件收敛, 而 $\sum\limits_{n=1}^\infty c_n$ 发散的例子}
取 $a_n = b_n = \dfrac{(-1)^{n+1}}{\sqrt{n}},$ 那么 $\sum\limits_{n=1}^\infty a_n$ 和 $\sum\limits_{n=1}^\infty b_n$ 都是条件收敛的, 而
\[
|c_n| = \left|\sum_{i+j=n+1} (-1)^{i+j+2} \dfrac{1}{\sqrt{i}\sqrt{j}}\right| = \left|(-1)^{n+1} \sum_{i+j=n+1} \dfrac{1}{\sqrt{i}\sqrt{j}}\right| \geqslant \sum_{i+j=n+1} \dfrac{1}{(i + j)/2} = \dfrac{2n}{n+1}.
\]
因此, $c_n \not\to 0,$ 故 $\sum\limits_{n=1}^\infty c_n$ 发散.
\end{block}

\end{frame}

%------------------------------------------------

\begin{frame}{绝对收敛级数的乘积}

\begin{question}
我们已经证明了
\begin{itemize}
\item 如果 $\sum\limits_{n=1}^\infty a_n = A, \sum\limits_{n=1}^\infty b_n = B$ 都收敛 ~ $\implies$ 正方形乘积 $ = A \cdot B$;
\item 如果 $\sum\limits_{n=1}^\infty a_n = A, \sum\limits_{n=1}^\infty b_n = B$ 都收敛, 且至少其中一个绝对收敛 ~ $\implies$ Cauchy 乘积 $ = A \cdot B$.
\end{itemize}
那么, 如果 $\sum\limits_{n=1}^\infty a_n = A, \sum\limits_{n=1}^\infty b_n = B$ 都绝对收敛, 那么会有更强的结论吗?
\end{question}

\pause

\begin{thm}
设 $\sum\limits_{n=1}^\infty a_n = A, \sum\limits_{n=1}^\infty b_n = B$ 都绝对收敛, 那么将 $a_k b_\ell$ 按照任意顺序排列成一个数列 $\{c_n\},$ 则 $\sum\limits_{n=1}^\infty c_n$ 也是绝对收敛的, 并且 $\sum\limits_{n=1}^\infty c_n = AB.$
\end{thm}

\end{frame}

%------------------------------------------------

\begin{frame}{绝对收敛级数的乘积}

\startproof[bg=yellow, fg=red](证明思路)
容易观察到, $c_n = a_{\sigma_1(n)} b_{\sigma_2(n)}$ 的前 $n$ 项和, 对应于 $a_n, b_n$ 排列成的无穷矩阵的区域, 可以用一个很大的正方形来包含, 这里的 $\sigma_1, \sigma_2$ 是 $a_n, b_n$ 的排列. 考虑 $|c_n|$ 的前 $N$ 项和
\[
\sum\limits_{n=1}^N |c_n| = \sum\limits_{n=1}^N |a_{\sigma_1(n)} b_{\sigma_2(n)}|,
\]
令 $M = \max\limits_{1 \leqslant n \leqslant N} \{\sigma_1(n), \sigma_2(n)\},$ 那么上式的绝对值可以被下面的式子控制:
\[
\sum\limits_{n=1}^N |c_n| \leqslant \left( \sum\limits_{n=1}^M |a_n| \right) \cdot \left( \sum\limits_{n=1}^M |b_n| \right) \leqslant \left( \sum\limits_{n=1}^\infty |a_n| \right) \cdot \left( \sum\limits_{n=1}^\infty |b_n| \right).
\]
由此可知正项级数 $\sum\limits_{n=1}^\infty |c_n|$ 收敛, 即 $\sum\limits_{n=1}^\infty c_n$ 绝对收敛.

\pause

绝对收敛级数的和不受项的排列影响, 将 $c_n$ 的次序调整成 Cauchy 乘积的次序, 则 $\sum\limits_{n=1}^\infty c_n = AB.$
\qed

\pause

\begin{question}
若 $\sum\limits_{n=1}^\infty a_n = A, \sum\limits_{n=1}^\infty b_n = B, \sum\limits_{n=1}^\infty c_n = C$ 都收敛, 那么是否必然有 $AB = C$? 或者是不一定成立的?
\end{question}

\end{frame}

%------------------------------------------------

\begin{frame}{级数乘积应用举例}

考虑 $e^x = \sum\limits_{n=0}^\infty \dfrac{x^n}{n!},$ $e^y = \sum\limits_{n=0}^\infty \dfrac{y^n}{n!}.$

首先, 由 d'Alembert 判别法, $\left.\dfrac{x^{n+1}}{(n+1)!} \right/ \dfrac{x^n}{n!} = \dfrac{x}{n+1} \to 0 ~ (n \to \infty),$ 知以上级数都是绝对收敛的.
\pause
考虑用 Cauchy 乘积的方式计算 $e^x \cdot e^y,$ 则
\begin{align*}
e^x \cdot e^y & = \left( \sum\limits_{n=0}^\infty \dfrac{x^n}{n!} \right) \cdot \left( \sum\limits_{n=0}^\infty \dfrac{y^n}{n!} \right) = \sum\limits_{n=0}^\infty \sum_{i+j=n} \dfrac{x^i}{i!} \cdot \dfrac{y^j}{j!} \cdot {\color{red} \dfrac{n!}{n!}} \\
& = \sum\limits_{n=0}^\infty \dfrac{1}{n!} \sum_{i+j=n} \dfrac{n!}{i!j!} x^i y^j = \sum\limits_{n=0}^\infty \dfrac{(x+y)^n}{n!} = e^{x+y}.
\end{align*}

\end{frame}

%------------------------------------------------

\section{级数求和的 Cesàro 方法}

%------------------------------------------------

\begin{frame}{级数求和的 Cesàro 方法的引入}

我们来扩展``可求和''级数的范围, 例如

\pause

\[\sum\limits_{n=1}^{\infty} (-1)^{n+1} = 1 - 1 + 1 - 1 + \cdots\]

\pause

这在通常意义下是一个发散的级数, 但我们很快会知道它是一个所谓的 Cesàro 意义下可和的级数.

\vspace{1em}
\pause

Cesàro 求和法, 实际上可以视作是某种``动态加权平均''的方法. 配合上 Tauber 型的定理, 可以在掌握数项级数某种意义上的前 $n$ 项均值变化趋势的前提下, 结合关于通项的某种``正则性'', 对级数本身的敛散性等性质做出判断.

\end{frame}

% %------------------------------------------------

\ifSiZheng

\begin{frame}{级数求和的 Cesàro 方法的引入}

\begin{columns}

\begin{column}{0.37\textwidth}
Tauber 型定理在数论中处理有关 Dirichlet 级数 (Riemann $\zeta$-函数, Dirichlet $L$-函数等) 有非常多的应用.

\pause
\vspace{0.4em}

我国著名数学家张寿武先生主要研究领域之一就是 $L$-函数. 同时, 他也培养了一大批国际上有影响力的青年数学家.
\end{column}
\begin{column}{0.6\textwidth}
\includegraphics[width=\textwidth]{images/shouwu-zhang.jpg}
\end{column}
\end{columns}

\end{frame}

\fi  % ifSiZheng

%------------------------------------------------

\begin{frame}{Cesàro 意义下可和的级数}

\begin{Def}[Cesàro 意义下可和的级数]
设 $s_n = \sum\limits_{k=1}^{n} a_k$ 是数项级数 $\sum\limits_{n=1}^{\infty} a_n$ 的前 $n$ 项和, 记 {\color{red}$\sigma_n = \dfrac{1}{n} \sum\limits_{k=1}^{n} s_k$}. 若 {\color{red}$\lim\limits_{n\to\infty} \sigma_n = A \in \mathbb{R}$}, 则称级数 $\sum\limits_{n=1}^{\infty} a_n$ 在 Cesàro 意义下可和, 或者更准确的说是在 $(c, 1)$ 的意义下可和, 并记为 $\sum\limits_{n=1}^{\infty} a_n = A~(c, 1).$
\end{Def}

\pause

$\sigma_n$ 实际上可以视作下面的加权平均:
\[
\begin{array}{r:cccccccccc}
  s_1 \phantom{++} & & a_1 \\
+ s_2 \phantom{++} & \phantom{+} + & a_1 & + & a_2 \\
% + s_3 \phantom{++} & \phantom{+} + & a_1 & + & a_2 & + & a_3 \\
\vdots \phantom{++} & & \vdots \\
+ s_n \phantom{++} & \phantom{+} + & a_1 & + & a_2 & + & a_3 & + & \cdots & + & a_n \\
\hdashline
\text{weight} \phantom{+} & & 1 & & \frac{n-1}{n} & & \frac{n-2}{n} & & \cdots & & \frac{1}{n}
\end{array}
\]

\end{frame}

%------------------------------------------------

\begin{frame}{Cesàro 意义下可和级数的例子}

% {\smaller[1]
考虑我们一开始提到的通常意义下发散的级数
\[\sum\limits_{n=1}^{\infty} a_n = \sum\limits_{n=1}^{\infty} (-1)^{n+1}  = 1 - 1 + 1 - 1 + \cdots.\]

\pause

容易算得
\[
s_n = \sum\limits_{k=1}^{n} a_k = \begin{cases}
1, & n = 2m - 1 \\
0, & n = 2m
\end{cases} \quad m = 1, 2, \dots
\]
进而有
\[
\sigma_n = \sum\limits_{k=1}^n s_k = \begin{cases}
\dfrac{m}{2m - 1}, & n = 2m - 1\\
\dfrac{1}{2}, & n = 2m
\end{cases} \quad m = 1, 2, \dots
\]
% }

\pause

$\sigma_n$ 关于 $n$ 有极限, 为 $\lim\limits \sigma_n = \dfrac{1}{2}.$ 于是这个级数在 $(c, 1)$ 意义下可和, 即有
\[\sum\limits_{n=1}^{\infty} a_n = \dfrac{1}{2}~(c, 1).\]

\end{frame}

%------------------------------------------------

\begin{frame}{Cesàro 意义下可和级数的性质}

一个自然的问题是: 对于通常意义下收敛的级数 $\sum\limits_{n=1}^{\infty} a_n = A \in \mathbb{R},$ 它在 $(c, 1)$ 意义下是可和的吗? 如果是, 那么它的和还是 $A$ 吗?

\pause
\vspace*{2em}

\begin{prop}
设级数 $\sum\limits_{n=1}^{\infty} a_n$ 在通常意义下收敛到有限实数 $A,$ 即 $\lim\limits_{n \to \infty} s_n = A,$ 那么级数 $\sum\limits_{n=1}^{\infty} a_n$ 是 $(c, 1)$ 可和的, 而且有 $\sum\limits_{n=1}^{\infty} a_n = A~(c, 1).$
\end{prop}

\end{frame}

%------------------------------------------------

\begin{frame}{Cesàro 意义下可和级数的性质}

\begin{proof}
若 $\sum\limits_{n=1}^{\infty} a_n$ 在通常意义下收敛到有限实数 $A,$ 那么
\[A = \lim\limits_{n \to \infty} s_n = \lim\limits_{n \to \infty} \dfrac{n\sigma_n - (n-1)\sigma_{n-1}}{n - (n-1)}, \]
于是根据 Stolz 公式 (Stolz–Cesàro 定理), 有
\[\lim\limits_{n \to \infty} \sigma_n = \lim\limits_{n \to \infty} \dfrac{n\sigma_n}{n} = \lim\limits_{n \to \infty} \dfrac{n\sigma_n - (n-1)\sigma_{n-1}}{n - (n-1)} = A.\]
这表明了 $\sum\limits_{n=1}^{\infty} a_n = A~(c, 1).$
\end{proof}

\end{frame}

%------------------------------------------------

\begin{frame}{$(c, r)$ 意义下可和的级数}

对于序列 $\{\sigma_k\},$ 可以考虑极限 $\lim\limits_{n\to\infty} \dfrac{1}{n} \sum\limits_{k=1}^n \sigma_k,$ 若此极限存在, 则称级数 $\sum\limits_{n=1}^{\infty} a_n$ 是 $(c, 2)$ 意义下可和的级数:
\[
\begin{array}{r:cccccccccc}
\sigma_1 \phantom{++} & & a_1 \\
+ \sigma_2 \phantom{++} & \phantom{+} + & a_1 & + & ~ {\color{red}\frac{1}{2}}a_2 \\
\vdots \phantom{++} & & \vdots \\
+ \sigma_n \phantom{++} & \phantom{+} + & a_1 & + & ~ {\color{red}\frac{n-1}{n}}a_2 & + & & \cdots & + & ~ {\color{red}\frac{1}{n}}a_n \\[4pt]
\hdashline
\text{weight} \phantom{+} & & 1 & & ~ {\color{red}\frac{1}{n}\sum\limits_{k=2}^n \frac{k-1}{k}} & & & \cdots & & ~ {\color{red}\frac{1}{n^2}}
\end{array}
\]

\pause
% \vspace*{0.4em}

由此, 可以归纳地定义 $(c, r)$ 意义下可和的级数, $r \geqslant 1$ 为正整数.
\[\sigma_n^{(r)} = \dfrac{1}{n} \sum\limits_{k=1}^n \sigma_n^{(r-1)} \text{ 关于 $n$ 收敛}, ~~ (\text{约定 } \sigma_n^{(1)} = \sigma_n, \sigma_n^{(0)} = s_n).\]

\end{frame}

%------------------------------------------------

\begin{frame}{$(c, r)$ 意义下可和的级数}

将通常意义下收敛的级数称作是 $(c, 0)$ 意义下可和的级数, 那么根据之前我们证明的命题以及举的例子, 我们有如下的集合 (空间) 的真包含的关系:
% \begin{multline*}
%   \{ (c, 0) \text{可和级数} \} \subsetneqq \{ (c, 1) \text{可和级数} \} \subsetneqq \cdots \\ \subsetneqq \{ (c, r) \text{可和级数} \} \subsetneqq \{ (c, r+1) \text{可和级数} \} \subsetneqq \cdots
% \end{multline*}

\[
\begin{array}{c}
\{ (c, 0) \text{可和级数} \} \\
\rotatebox[origin=c]{-90}{$\subsetneqq$} \\
\{ (c, 1) \text{可和级数} \} \\
\rotatebox[origin=c]{-90}{$\subsetneqq$} \\
\vdots \\
\{ (c, r) \text{可和级数} \} \\
\rotatebox[origin=c]{-90}{$\subsetneqq$} \\
\{ (c, r+1) \text{可和级数} \} \\
\rotatebox[origin=c]{-90}{$\subsetneqq$} \\
\vdots
\end{array}
\]

\end{frame}

%------------------------------------------------

\section[Tauber 型定理]{Tauber 型定理}

%------------------------------------------------

\begin{frame}{Tauber 型定理}

反过来, 我们有一些结论, 对 $(c, r+1)$ 可和的级数, 加上一些额外的限制, 能得到它落在子集, $\{ (c, r) \text{可和级数} \}$ 中. 这就是所谓的 Tauber 型的定理.

\pause

\vspace*{1em}

\begin{thm}[Tauber]
若级数 $\sum\limits_{n=1}^{\infty} a_n = A~(c, 1),$ 且当 $n\to\infty$ 时 $a_n = \mathrm{o} \left( \dfrac{1}{n} \right),$ 那么该级数在通常意义下可和, 且 $\sum\limits_{n=1}^{\infty} a_n = A.$
\end{thm}

\end{frame}

%------------------------------------------------

\begin{frame}{Tauber 型定理证明}

我们考察 $s_n - \sigma_n,$
\pause
有
\begin{align*}
s_n - \sigma_n & = \sum\limits_{k=1}^n a_k - \sum\limits_{k=1}^n \left( 1 - \dfrac{k-1}{n} \right) a_k \\
& = \dfrac{\sum\limits_{k=1}^n (k-1)a_k}{n}.
\end{align*}
\pause
由于 $n \to \infty$ 时有 $a_n = \mathrm{o}\left( \dfrac{1}{n} \right),$ 那么 $\lim\limits_{n \to \infty} n a_n = 0,$ 这也等价于
\[\lim\limits_{n \to \infty} (n - 1) a_n = 0.\]

\end{frame}

%------------------------------------------------

\begin{frame}{Tauber 型定理证明}

于是对 $\dfrac{\sum\limits_{k=1}^n (k-1)a_k}{n}$ 使用 Stolz 公式有:
\begin{align*}
\lim\limits_{n \to \infty} (s_n - \sigma_n) & = \lim\limits_{n \to \infty} \dfrac{\sum\limits_{k=1}^n (k-1)a_k}{n} \\
& = \lim\limits_{n \to \infty} \dfrac{\sum\limits_{k=1}^n (k-1)a_k - \sum\limits_{k=1}^{n-1} (k-1)a_k}{n - (n-1)} \\
& = \lim\limits_{n \to \infty} \dfrac{(n-1)a_n}{1} = 0.
\end{align*}
这表明 $\sum\limits_{n=1}^{\infty} a_n = \lim\limits_{n \to \infty} s_n = \lim\limits_{n \to \infty} \sigma_n = A.$ \qed

\end{frame}

%------------------------------------------------

\begin{frame}{更强的 Tauber 型定理}

实际上, 我们可以减弱 Tauber 定理中关于 $a_n = \mathrm{o} \left( \dfrac{1}{n} \right)$ 的假设, 得到一个更强的结果:

\begin{thm}[Hardy]
若级数 $\sum\limits_{n=1}^{\infty} a_n = A~(c, 1),$ 且当 $n\to\infty$ 时 $a_n = \mathrm{O} \left( \dfrac{1}{n} \right),$ 那么该级数在通常意义下可和, 且 $\sum\limits_{n=1}^{\infty} a_n = A.$
\end{thm}

\vspace*{1em}

注意: 这里关于 $a_n$ 趋于 $0$ 速度的假设, 我们从 ${\color{red} \lim\limits_{n\to\infty} n a_n = 0},$ 减弱为了: 存在正实数 $M,$ 使得对任意自然数 $n$ 有 ${\color{red} \rvert n a_n \rvert \leqslant M}.$

\end{frame}

%------------------------------------------------

\begin{frame}{增强的 Tauber 型定理的分析与证明}

除了已知条件: 当 $n\to\infty$ 时 $a_n = \mathrm{O} \left( \dfrac{1}{n} \right),$ 即存在正实数 $M,$ 使得对任意自然数 $n$ 有 ${\color{red} \rvert n a_n \rvert \leqslant M}$ 之外, 我们所拥有的条件是 $\sum\limits_{n=1}^{\infty} a_n = A~(c, 1),$ 即 $\lim\limits_{n\to\infty} \sigma_n = A,$ 这等价于 $\forall~\varepsilon > 0,$ 存在自然数 $N$ 使得 ${\color{red} \forall~n > m > N, \lvert \sigma_n - \sigma_m \vert < \varepsilon}.$

\pause
\vspace*{0.4em}

之前证明 Tauber 定理的时候, 我们导出了 $s_n$ 与 $\sigma_n$ 的关系
\[s_n = \sigma_n + \sum\limits_{k = 1}^n \dfrac{(k-1)a_k}{n}.\]
由于我们只有 $\rvert n a_n \rvert \leqslant M,$ 所以 Stolz 公式不能再用了, 得另想办法.

\pause
\vspace*{0.4em}

\begin{attnblock}
很多判别法、计算公式, 例如 L'Hopital 法则, Stolz 公式, 包括后面要学的一致收敛的函数项级数的逐项微分、积分的结论, 相关条件都是一种充分性条件, 远 \textbf{\color{red}不是必要条件}.
\end{attnblock}

\end{frame}

%------------------------------------------------

\begin{frame}{增强的 Tauber 型定理的分析与证明}

想法: 将 $s_n - \sigma_n$ 与 $\sigma_n - \sigma_m$ 联系上, 并用它控制余项的大小.

\only<2>{
假设 $n > m,$ 那么有
\begin{align*}
\sigma_n - \sigma_m & = \dfrac{1}{n} \sum\limits_{k=1}^n s_k - \dfrac{1}{m} \sum\limits_{k=1}^{{\color{red}m}} s_k \\
& = \dfrac{1}{n} \sum\limits_{k=1}^n s_k - \dfrac{1}{m} \sum\limits_{k=1}^{{\color{red}n}} s_k + \dfrac{1}{m} \sum\limits_{{\color{red}k=m+1}}^{n} s_k \\
& = \dfrac{m - n}{nm} \sum\limits_{k=1}^n s_k + \dfrac{1}{m} \sum\limits_{k=m+1}^{n} s_k.
\end{align*}

另一方面, 有
\[s_n - \sigma_n = s_n - \dfrac{1}{n} \sum\limits_{k = 1}^n s_k.\]
}

\only<3>{假设 $n > m,$ 那么有
\begin{align*}
\sigma_n - \sigma_m & = \dfrac{1}{n} \sum\limits_{k=1}^n s_k - \dfrac{1}{m} \sum\limits_{k=1}^m s_k \\
& = \dfrac{1}{n} \sum\limits_{k=1}^n s_k - \dfrac{1}{m} \sum\limits_{k=1}^{n} s_k + \dfrac{1}{m} \sum\limits_{k=m+1}^{n} s_k \\
& = \dfrac{m - n}{nm} {\color{red}
\sum\limits_{k=1}^n s_k} + \dfrac{1}{m} \sum\limits_{k=m+1}^{n} s_k.
\end{align*}

另一方面, 有
\[s_n - \sigma_n = s_n - \dfrac{1}{n} {\color{red} \sum\limits_{k = 1}^n s_k}.\]
}

\end{frame}

%------------------------------------------------

\begin{frame}{增强的 Tauber 型定理的分析与证明}

两式通分相减, 有
\begin{align*}
& (s_n - \sigma_n) - \dfrac{m}{n - m}(\sigma_n - \sigma_m) \\
& = s_n - {\color{red} \cancel{\dfrac{1}{n} \sum\limits_{k = 1}^n s_k}} + {\color{red} \cancel{\dfrac{1}{n} \sum\limits_{k=1}^n s_k}} - \dfrac{1}{n-m} \sum\limits_{k=m+1}^{n} s_k \\
& = s_n - \dfrac{1}{n-m} \sum\limits_{k=m+1}^{n} s_k = \dfrac{1}{n-m} \sum\limits_{k=m+1}^{n} (s_n - s_k).
\end{align*}

\pause
\vspace*{1em}

整理一下, 有
\[s_n - \sigma_n = \underbrace{\color{cyan}\dfrac{m}{n - m}(\sigma_n - \sigma_m)}_{\text{\ding{172}}} + \underbrace{\color{magenta} \dfrac{1}{n-m} \sum\limits_{k=m+1}^{n} (s_n - s_k)}_{\text{\ding{173}}}.\]

\end{frame}

%------------------------------------------------

\begin{frame}{增强的 Tauber 型定理的分析与证明}

对于 $\text{\ding{173}} = \dfrac{1}{n-m} \sum\limits_{k=m+1}^{n} (s_n - s_k)$ 和式中的每一项, 我们有估计
{\smaller
\begin{multline*}
\lvert s_n - s_k \rvert = \lvert a_{k+1} + \cdots + a_n \rvert \\
\leqslant \dfrac{M}{k+1} + \cdots \dfrac{M}{n} \leqslant \dfrac{(n-k)M}{k+1} \leqslant \dfrac{(n-m-1)M}{m+1}.
\end{multline*}
}
于是
{\smaller
\[
\left\lvert \dfrac{1}{n-m} \sum\limits_{k=m+1}^{n} (s_n - s_k) \right\rvert \leqslant \dfrac{1}{n-m} (n-m) \dfrac{n-m-1}{m+1} M < \dfrac{n-m}{m} M.
\]
}

\pause

我们希望有 $\dfrac{n-m}{m} < \dfrac{\varepsilon}{2M},$ 即 $m > \dfrac{n}{1 + \frac{\varepsilon}{2M}}.$ 取 $n$ 足够大, 使得 $n - \dfrac{n}{1 + \frac{\varepsilon}{2M}} > 2,$ 或者等价地, 取 $\color{red} n > \dfrac{2(2M + \varepsilon)}{\varepsilon}$, 即可确保能取到整数 $\color{red} m \in \left[ \dfrac{n}{1 + \frac{\varepsilon}{2M}}, n \right]$,
即有 $\lvert \text{\ding{173}} \rvert < \dfrac{\varepsilon}{2}.$

\end{frame}

%------------------------------------------------

\begin{frame}{增强的 Tauber 型定理的分析与证明}

对于 $\text{\ding{172}} = \dfrac{m}{n - m}(\sigma_n - \sigma_m),$ 由于已有 $\dfrac{n-m}{m} < \dfrac{\varepsilon}{2M},$ 即 $\dfrac{m}{n-m} > \dfrac{2M}{\varepsilon},$ 进一步要求 $\dfrac{2M}{\varepsilon} < \dfrac{m}{n-m} < \dfrac{4M}{\varepsilon},$ 即 $m < \dfrac{n}{1 + \frac{\varepsilon}{4M}}.$
\pause
这样的整数 $m$ 总是可以取到的, 只要保证 $\dfrac{n}{1 + \frac{\varepsilon}{4M}} - \dfrac{n}{1 + \frac{\varepsilon}{2M}} > 2$ 即可, 即 $\color{red} n > \dfrac{8M}{\varepsilon}\left(1 + \frac{\varepsilon}{4M}\right)\left(1 + \frac{\varepsilon}{2M}\right)$.
\pause
于是
\[\lvert \text{\ding{172}} \rvert = \left\lvert \dfrac{m}{n - m}(\sigma_n - \sigma_m) \right\rvert \leqslant \dfrac{4M}{\varepsilon} \left\lvert \sigma_n - \sigma_m \right\rvert.\]
若取 $N$ 充分大, 使得 $\forall~n > m > N$ 都有 $\left\lvert \sigma_n - \sigma_m \right\rvert < {\color{red} \dfrac{\varepsilon^2}{8M}},$ 那么
\[\left\lvert s_n - \sigma_n \right\rvert \leqslant \lvert \text{\ding{172}} \rvert + \lvert \text{\ding{173}} \rvert < \varepsilon.\]

\pause
% \vspace{1em}

注意这里的关键是 $m, n$ 的取法, 满足 $n$ 充分大, 且
\[\dfrac{n}{1 + \frac{\varepsilon}{2M}} < m < \dfrac{n}{1 + \frac{\varepsilon}{4M}} < n.\]

% 将上述思考、分析过程整理、调整顺序之后, 即是完整的证明.

\end{frame}

%------------------------------------------------

\begin{frame}{Tauber 型定理}

以上我们证明了两个所谓的 Tauber 型的定理, 即在某些附加的``正则性''条件下, 根据考察的量的某些均值的性质, 对这个量本身的性质作一些判断.

\pause
\vspace{1em}

这里, 我们附加了
\[a_n = s_n - s_{n-1} = \mathrm{o}\left(\dfrac{1}{n}\right) \text{ 或 } \mathrm{O}\left(\dfrac{1}{n}\right)\]
的``正则性''条件, 根据 $s_n$ 的均值 $\sigma_n$ 的收敛性, 得到了 $s_n$ 本身的收敛性.

\end{frame}

% %------------------------------------------------

% \section{总结}

% %------------------------------------------------

% \begin{frame}{总结}

% {\larger[2] \color{red} \ding{43}} 级数的Cauchy收敛原理 (充要条件):
% \[\left\lvert x_{n+1} + x_{n+2} + \cdots + x_m \right\rvert = \left\lvert \sum\limits_{k=n+1}^m x_k \right\rvert < \varepsilon\]

% \pause

% {\larger[2] \color{red} \ding{43}} Leibniz判别法: Leibniz级数都收敛.
% \[\text{Leibniz级数:} \sum\limits_{n=1}^{\infty} (-1)^{n+1} u_n, ~~ 0 < u_n \text{ 单调减小, 且 } u_n \to 0.\]

% \pause

% {\larger[2] \color{red} \ding{43}} A-D判别法 ($\sum\limits_{n=1}^{\infty} a_n b_n$收敛的充分条件):
% \begin{align*}
% \text{Abel判别法:} & \sum\limits_{n=1}^{\infty} b_n ~ \text{收敛}, ~ a_n ~ \text{单调有界} \\
% \text{Dirichlet判别法:} & \sum\limits_{n=1}^{N} b_n ~ \text{关于$N$有界}, ~ a_n ~ \text{单调趋于$0$}
% \end{align*}

% \end{frame}

%------------------------------------------------

%------------------------------------------------
% % closing page
% \begin{frame}
% \HUGE{\centerline{\bfseries {\color{gray} The} {\color{zkblue} End}}}

% \pause

% \vspace{1.2em}

% \Large{\centerline{\bfseries {\color{gray}Thank} \quad {\color{zkblue} You}}}

% \end{frame}

%----------------------------------------------------------------------------------------

\end{document}
